\documentclass{article}
\usepackage{ctex, amsmath, amssymb, array, booktabs, geometry,enumitem,tasks}
\geometry{a4paper, margin=1.5cm}
\usepackage{titling}\setlength{\droptitle}{-2cm}

\usepackage{enumitem}
\setlist[itemize]{itemsep=0pt, parsep=0pt, topsep=3pt}
\setlist[enumerate]{itemsep=0pt, parsep=0pt, topsep=3pt}

\author{学号：\underline{\hspace{4cm}} 姓名：\underline{\hspace{4cm}} }
\title{偏微分方程与 Black-Scholes 期权定价公式}
\date{2026年10月21日}
\begin{document}
\maketitle

\begin{enumerate}[label=\textbf{\arabic*.}]

% \section*{一、扩散方程与 Delta 函数}

\item  {\bf{【验证函数是否为扩散方程的解】}} %题目 1. 

扩散方程（热传导方程）为：
\(
\frac{\partial u}{\partial \tau} = \frac{\partial^2 u}{\partial x^2}, 
\)
已知函数：
\(
u(x, \tau) = \frac{1}{2\sqrt{\pi\tau}} e^{-\frac{x^2}{4\tau}}, \quad -\infty < x < +\infty, \; \tau > 0. 
\)

\begin{enumerate}[label=(\arabic*)]
    \item 求 \(\frac{\partial u}{\partial \tau}\)；
    \item 求 \(\frac{\partial^2 u}{\partial x^2}\)；
    \item 验证 \(u(x, \tau)\) 是否满足扩散方程。
\end{enumerate}



\item  {\bf{【由相似解方法求解半无限区间扩散方程】}} %题目 2. 

考虑半无限区间上的扩散方程：
\(
\frac{\partial u}{\partial \tau} = \frac{\partial^2 u}{\partial x^2}, \quad x > 0, \tau > 0. 
\)

初始条件：
\(
u(x, 0) = 1, 
\)
边界条件：
\(
u(0, \tau) = 0, \quad u \to 0 \; (x \to +\infty). 
\)

设解具有相似形式 \(u(x, \tau) = U(\xi)\)，其中 \(\xi = \frac{x}{\sqrt{\tau}}\). 

\begin{enumerate}[label=(\arabic*)]
    \item 利用链式法则求 \(\frac{\partial u}{\partial \tau}\) 和 \(\frac{\partial^2 u}{\partial x^2}\)；
    \item 将偏微分方程化为常微分方程；
    \item 求解该常微分方程并写出 \(u(x, \tau)\) 的表达式。
\end{enumerate}



\item  {\bf{【由 delta 函数初始条件推导扩散方程的基本解】}} %题目 3. 

设扩散方程：
\(
\frac{\partial u}{\partial \tau} = \frac{\partial^2 u}{\partial x^2}
\)
的初始条件为 delta 函数：
\(
u(x, 0) = \delta(x). 
\)

设解具有相似形式 \(u(x, \tau) = \tau^{-1/2} U(\xi)\)，其中 \(\xi = \frac{x}{\sqrt{\tau}}\). 

\begin{enumerate}[label=(\arabic*)]
    \item 求 \(\frac{\partial u}{\partial \tau}\) 和 \(\frac{\partial^2 u}{\partial x^2}\)；
    \item 化为常微分方程并求解；
    \item 利用归一化条件 \(\int_{-\infty}^{+\infty} u(x, \tau) dx = 1\) 确定常数，写出基本解。
\end{enumerate}



\item  {\bf{【计算 delta 函数与光滑函数的积分】}} %题目 4. 

已知 delta 函数的性质：
\(\displaystyle 
\int_{-\infty}^{+\infty} \delta(x - x_0) \Phi(x) dx = \Phi(x_0). 
\)

计算下列积分：

\begin{enumerate}[label=(\arabic*)]
    \item \(\int_{-\infty}^{+\infty} \delta(x - 3) (x^2 + 2x) dx\)；
    \item \(\int_{0}^{5} \delta(x - 2) e^{x} dx\)；
    \item \(\int_{-\infty}^{+\infty} \delta(x) \cos(x) dx\)。
\end{enumerate}



\item  {\bf{【由 Delta 函数表示不连续函数的导数】}} %题目 5. 

某人的财富 \(M(t)\) 满足：
\(\displaystyle 
M(t) = \begin{cases} 0, & 0 < t < t_0, \\ D_0, & t \geqslant t_0. \end{cases} 
\)

\begin{enumerate}[label=(\arabic*)]
    \item 写出 \(M(t)\) 的导数 \(\frac{dM}{dt}\)，用 delta 函数表示；
    \item 若 \(D_0 = 1000\)，\(t_0 = 5\)，计算 \(\int_0^{10} \frac{dM}{dt} dt\). 
\end{enumerate}


% \section*{二、Black-Scholes 微分方程推导}

\item  {\bf{【构造无风险资产组合并推导 Black-Scholes 微分方程】}} %题目 6. 

设股票价格 \(S\) 遵循几何布朗运动：
\(
dS = \mu S dt + \sigma S dz. 
\)

设期权价格为 \(V(S, t)\)，由 Itô 引理：
\(\displaystyle 
dV = \left( \frac{\partial V}{\partial S} \mu S + \frac{\partial V}{\partial t} + \frac{1}{2} \frac{\partial^2 V}{\partial S^2} \sigma^2 S^2 \right) dt + \frac{\partial V}{\partial S} \sigma S dz. 
\)

构造投资组合：
\(
\pi = V - \delta S. 
\)

\begin{enumerate}[label=(\arabic*)]
    \item 求 \(d\pi = dV - \delta dS\)；
    \item 选择 \(\delta = \frac{\partial V}{\partial S}\)，消去 \(dz\) 项；
    \item 由无套利条件 \(d\pi = r\pi dt\)，推导 Black-Scholes 微分方程。
\end{enumerate}




% \section*{三、Black-Scholes 欧式期权定价公式计算}

\item  {\bf{【用 Black-Scholes 公式计算欧式看涨期权价格】}} %题目 7. 

当前股票价格 \(S = 50\) 元，欧式看涨期权的执行价格 \(E = 45\) 元，无风险年利率 \(r = 6\%\)， \\ 
期权有效期 \(T - t = 0.25\) 年，股票价格波动率 \(\sigma = 0.2\)。

已知 Black-Scholes 公式：
\(
C(S, t) = SN(d_1) - Ee^{-r(T-t)} N(d_2),  
\)
其中 
\[
d_1 = \frac{\ln(S/E) + \left( r + \frac{1}{2}\sigma^2 \right)(T-t)}{\sigma \sqrt{T-t}}, \quad 
d_2 = d_1 - \sigma \sqrt{T-t}. 
\]

\begin{enumerate}[label=(\arabic*)]
    \item 计算 \(d_1\) 和 \(d_2\)；
    \item 查表得 \(N(d_1) \approx 0.7422\)，\(N(d_2) \approx 0.6651\)，计算期权价格 \(C\)。
\end{enumerate}



\item  {\bf{【用 Black-Scholes 公式计算欧式看跌期权价格】}} %题目 8. 

沿用题目 7 的数据：\(S = 50\) 元，\(E = 45\) 元，\(r = 6\%\)，\(T - t = 0.25\) 年，\(\sigma = 0.2\)。

由 Black-Scholes 公式，欧式看跌期权价格为：
\(
P(S, t) = Ee^{-r(T-t)} N(-d_2) - SN(-d_1). 
\)

已知 \(N(-d_1) = 1 - N(d_1)\)，\(N(-d_2) = 1 - N(d_2)\)，且 \(N(d_1) \approx 0.7422\)，\(N(d_2) \approx 0.6651\)。

\begin{enumerate}[label=(\arabic*)]
    \item 计算 \(N(-d_1)\) 和 \(N(-d_2)\)；
    \item 计算看跌期权价格 \(P\)；
    \item 用平价公式验证：\(P = C - S + Ee^{-r(T-t)}\)。
\end{enumerate}



\item  {\bf{【求 Black-Scholes 公式中的 \(d_1\) 和 \(d_2\)】}} %题目 9. 

已知某欧式看涨期权的参数如下：
股票价格 \(S = 100\)；
执行价格 \(E = 95\)；
无风险利率 \(r = 10\%\)；\\ 
到期时间 \(T - t = 0.25\) 年；
波动率 \(\sigma = 0.5\)。

\begin{enumerate}[label=(\arabic*)]
    \item 计算 \(\ln(S/E)\)；
    \item 计算 \(d_1\)；
    \item 计算 \(d_2\)。
\end{enumerate}



\item  {\bf{【分析波动率对期权价格的影响】}} %题目 10. 

某欧式看涨期权的参数如下：
\(S = 100\)，\(E = 95\)，\(r = 0.1\)，\(T - t = 0.25\)。

已知当 \(\sigma = 0.5\) 时，\(d_1 \approx 0.4302\)，\(d_2 \approx 0.1802\)，\(N(d_1) \approx 0.6664\)，\(N(d_2) \approx 0.5714\)。

当 \(\sigma = 0.6\) 时，\(d_1 \approx 0.4043\)，\(d_2 \approx 0.1043\)，\(N(d_1) \approx 0.6570\)，\(N(d_2) \approx 0.5415\)。

\begin{enumerate}[label=(\arabic*)]
    \item 分别计算两种波动率下的看涨期权价格；
    \item 说明波动率增大对看涨期权价格的影响。
\end{enumerate}


\item  {\bf{【解释 Black-Scholes 微分方程的含义】}} %题目 11. 

回答下列概念问题：

\begin{enumerate}[label=(\arabic*)]
    \item 写出 Black-Scholes 期权微分方程；
    \item 解释方程中各项的金融含义；
    \item 为什么 Black-Scholes 微分方程中不包含股票的期望收益率 \(\mu\)？这一事实有什么重要含义？
\end{enumerate}




\item  {\bf{【利用风险中性定价原理计算期权价格】}} %题目 12. 

设某不支付红利股票当前价格 \(S = 50\) 元，欧式看涨期权执行价格 \(E = 50\) 元，无风险利率 \(r = 10\%\)， \\ 
到期时间 \(T - t = 0.5\) 年，波动率 \(\sigma = 0.3\)。

已知 Black-Scholes 公式：
\(
C = SN(d_1) - Ee^{-r(T-t)}N(d_2)。 
\)

计算得 \(d_1 \approx 0.3005\)，\(d_2 \approx 0.0884\)，\(N(d_1) \approx 0.6181\)，\(N(d_2) \approx 0.5352\)。

\begin{enumerate}[label=(\arabic*)]
    \item 计算该欧式看涨期权的价格；
    \item 用平价公式计算相应的欧式看跌期权价格；
    \item 解释风险中性定价原理的核心思想。
\end{enumerate}


\end{enumerate}

\end{document}